\documentclass[11pt,a4paper]{article}

\usepackage[utf8]{inputenc}
\usepackage[T1]{fontenc}
\usepackage[spanish]{babel}
\usepackage{amsmath}
\usepackage{amsfonts}
\usepackage{amssymb}
\usepackage{graphicx}
\usepackage{tabularx}
\usepackage{booktabs}
\usepackage{hyperref}
\usepackage{geometry}
\usepackage{listings}
\usepackage{xcolor}
\usepackage{caption}
\usepackage{subcaption}
\usepackage{titlesec}
\usepackage{parskip}
\usepackage{float}

% Configuración de títulos de párrafos
\titleformat{\paragraph}[hang]
{\normalfont\normalsize\bfseries}{\theparagraph}{1em}{}
\titlespacing*{\paragraph}
{0pt}{3.25ex plus 1ex minus .2ex}{1em}

\geometry{top=3cm, bottom=3cm, left=2.5cm, right=2.5cm}

\hypersetup{
    colorlinks=true,
    linkcolor=black,
    filecolor=magenta,      
    urlcolor=cyan,
}

\definecolor{codegreen}{rgb}{0,0.6,0}
\definecolor{codegray}{rgb}{0.5,0.5,0.5}
\definecolor{codepurple}{rgb}{0.58,0,0.82}
\definecolor{backcolour}{rgb}{0.95,0.95,0.92}

\lstdefinestyle{mystyle}{
    backgroundcolor=\color{backcolour},   
    commentstyle=\color{codegreen},
    keywordstyle=\color{magenta},
    numberstyle=\tiny\color{codegray},
    stringstyle=\color{codepurple},
    basicstyle=\ttfamily\small,
    breakatwhitespace=false,         
    breaklines=true,                 
    captionpos=b,                    
    keepspaces=true,                 
    numbers=left,                    
    numbersep=5pt,                  
    showspaces=false,                
    showstringspaces=false,
    showtabs=false,                  
    tabsize=2
}

\lstset{style=mystyle}

\begin{document}
\begin{titlepage}
    \begin{center}
        \vspace*{1cm}

        \Huge
        \textbf{Laboratorio N°1}

        \vspace{0.5cm}
        \LARGE
        Muestreo con FRDM-MCXN947

        \vspace{1.5cm}

        \vfill

        \Large
        \textbf{Autor:} Simón Saillen

        \Large
        \textbf{Profesores:} Gustavo Parlanti, Roberto Rossi, German Molina, Pablo Ferreira

        \vspace{0.5cm}

        \Large
        \textbf{Procesamiento Digital de Señales}\\
        \vspace{0.2cm}
        \textbf{FCEFyN -- UNC}\\
        \vspace{0.2cm}
        2026

    \end{center}
\end{titlepage}
\newpage

\section{Introducción}

En este laboratorio se implementará un sistema de adquisición de señales
analógicas utilizando la placa FRDM-MCXN947. Se explorarán diferentes
velocidades de muestreo y se analizará el efecto de estas velocidades en la
calidad de la señal adquirida.

\section{Consigna}

Realizar la lectura de una señal analogica a distintas velocidades de muestreo
en KiloSamples/segundo: \texttt{8 kS/s}, \texttt{16 kS/s}, \texttt{22 kS/s},
\texttt{44 kS/s} y \texttt{48 kS/s}.

Los cambios de las velocidades de muestreo serán realizados con una de las
teclas de la placa, en forma de un buffer circular. Cada velocidad de muestreo
se indicara a traves de un color RGB del LED integrado.

Con otra tecla de la placa se habilitara la adquisicion y paro de muestras
(\textbf{RUN} / \textbf{STOP}).

Los valores adquiridos seran almacenados en memoria en un buffer circular de
512 muestras del tipo \texttt{q15} (fraccional de 15 bits) y a su vez seran
enviados a traves del DAC de 12 bits.

\section{Desarrollo}

El desarrollo se encuentra separado en distintas secciones, cada una de ellas
explicando la configuracion de los distintos modulos de la placa.

\subsection{Configuración de Clocks}

Principalmente se deben configurar los clocks de la placa, ya que todos los
modulos de la misma dependen de esta configuracion. Para ello se utilizara
dentro de ConfigTools el modulo \texttt{CLOCK} del SDK de NXP, el cual permite
configurar los distintos clocks de la placa.

\begin{table}[H]
    \centering
    \begin{tabular}{|c|c|c|c|}
        \hline
        \textbf{Clock Component} & \textbf{Frequency} & \textbf{Clock Source} & \textbf{Description}        \\
        \hline
        MAIN\_clock              & 150 MHz            & PLL0                  & Clock principal de la placa \\
        ADC0\_clock              & 48 MHz             & Fast IRC              & Clock del modulo ADC        \\
        DAC0\_clock              & 48 MHz             & Fast IRC              & Clock del modulo DAC        \\
        CTIMER0\_clock           & 150 MHz            & PLL0                  & Clock del modulo Timer      \\
        FLEXXCOM4\_clock         & 12 MHz             & Slow IRC              & Clock del modulo UART       \\
        \hline
    \end{tabular}
    \caption{Configuración de clocks usando ConfigTools}
\end{table}

El modulo UART se configuro para enviar mensajes de debugging a la PC, por lo
que se utilizo una pre configuracion con un baudrate de \texttt{115200 bps} y
un clock de 12 MHz.

\subsection{Modulo Low-Power ADC (LPADC0)}

El modulo ADC de la placa FRDM-MCXN947 es un conversor de 12 bits de baja
potencia con un rango de voltaje de refencia de 0 a 3.3V. El mismo se configuro
con los siguientes parametros:

\begin{table}[H]
    \centering
    \begin{tabular}{|c|c|c|}
        \hline
        \textbf{Campo}          & \textbf{Valor} & \textbf{Descripcion}                                \\
        \hline
        Pre-enable ADC          & Enabled        & Enciende el circuito analogico con anterioridad     \\
        Autocalibration         & Enabled        & Calibra al iniciar para compensar desviaciones      \\
        Doze mode               & Enabled        & Permite que funcione en modo de bajo consumo        \\
        FIFO 0 watermark        & 1              & Genera interrupcion cuando hay 1 muestra en la FIFO \\
        Enable interrupt vector & Enabled        & Habilita la interrupcion del ADC                    \\
        Interrupt sources       & FIFO 0         & La fuente de interrupcion es la FIFO 0 watermark    \\
        Interrupt priority      & 0              & Prioridad de la interrupcion del ADC                \\
        \hline
    \end{tabular}
    \caption{Configuración General del ADC en ConfigTools/Peripherals}
\end{table}

Al generar la interrupción, el CPU ejecuta la rutina de servicio de
interrupción (ISR) del ADC, la cual se encarga de leer la muestra de la FIFO y
almacenarla en el buffer circular de 512 muestras y pasarla tambien al DAC para
que la muestre en su salida.

\begin{table}[H]
    \centering
    \begin{tabular}{|c|c|c|}
        \hline
        \textbf{Campo}      & \textbf{Valor}                  & \textbf{Descripcion}        \\
        \hline
        Command ID          & start\_conversion               & Identificador del comando   \\
        Command number      & 1                               & Numero de comando           \\
        Channel sample mode & Single end mode, using A        & Modo de muestreo            \\
        Channel number      & ADC0\_A0                        & Pin de entrada del ADC      \\
        Resolution mode     & High resolution (16 bits)       & Resolucion de la conversion \\
        Loop count          & 0                               & Ejecuta el comando una vez  \\
        Sample time mode    & 3 ADCK cycles total sample time & Tiempo de muestreo          \\
        \hline
    \end{tabular}
    \caption{Conversion Command del ADC en ConfigTools/Peripherals}
\end{table}

\begin{table}[H]
    \centering
    \begin{tabular}{|c|c|c|}
        \hline
        \textbf{Campo}         & \textbf{Valor}  & \textbf{Descripcion}                  \\
        \hline
        Trigger ID             & timer0\_trigger & Identificador del trigger             \\
        Trigger number         & 0               & Numero de trigger                     \\
        Trigger command number & 1               & Numero de comando a ejecutar          \\
        High priority          & Enabled         & Ejecuta el comando con prioridad alta \\
        Select channel A FIFO  & FIFO 0          & Habilita la FIFO A del ADC            \\
        Delay value            & 0               & No hay delay entre trigger y comando  \\
        \hline
    \end{tabular}
    \caption{Trigger Command del ADC en ConfigTools/Peripherals}
\end{table}

\subsubsection{q15: Fraccional de 15 bits}

El formato \texttt{Q15} (o \texttt{Q1.15} o \texttt{S(16,15)}) es una notación
de punto fijo con signo que usa 1 bit para la parte entera y 15 bits para la
parte fraccionaria. Este formato permite representar números en el rango de
$-1$ a $(1 - 2^{-15}) \approx 1$, con una resolución de $2^{-15} \approx
    0{,}00003051$.

El ADC de la placa es un conversor de 16 bits (en high resolution mode), dicho
conversor no puede interpretar valores negativos (para solucionar esto,
generalmente la señal de entrada es movida al rango positivo por hardware
insertando un offset de voltaje continuo), para poder representar los valores
de la señal de entrada en el formato q15 (que si contempla valores negativos)
se debe restar a cada muestra el valor de $2^{15} = 32768$ (quitarle el
offset). Asi se pasa de un rango de [\texttt{0}, \texttt{0{x}FFFF}] a un rango
de [\texttt{0{x}8000}, \texttt{0{x}7FFF}], que es el rango de valores que puede
representar el formato q15.

\subsection{Modulo Timer (CTIMER0)}

El modulo CTIMER0 es un timer de 32 bits que se utiliza para generar el trigger
de conversión del ADC.\@ Se configuró desde Peripherals en ConfigTools con los
siguientes parámetros:

\begin{table}[H]
    \centering
    \begin{tabular}{|c|c|c|}
        \hline
        \textbf{Campo}          & \textbf{Valor}           & \textbf{Descripcion}               \\
        \hline
        Timer mode              & Timer (bus clock source) & Modo de funcionamiento             \\
        Prescaler value         & 1                        & Frecuencia de clock del timer      \\
        Match channel           & 0                        & Canal de match del timer           \\
        Match value             & 18750 ($\approx 8kHz$)   & Valor de match del timer           \\
        Reset on match          & Enabled                  & Resetea el timer al hacer match    \\
        Match interrupt request & Enabled                  & Genera interrupcion al hacer match \\
        Interrupt priority      & 1                        & Prioridad de la interrupcion       \\
        \hline
    \end{tabular}
    \caption{Configuración del Timer en ConfigTools/Peripherals}
\end{table}

El match value inicial corresponde a la velocidad de muestreo de \texttt{8
    kS/s}, posteriorment con el uso de los switch este valor es cambiado para
cumplir con las velocidades de muestreo requeridas.

\subsection{Modulo Low-Power DAC (LPDAC0)}

El modulo DAC de la placa es un conversor digital a analógico de 12 bits con un
rango de voltaje de salida de 0 a 3.3V. El mismo se configuró desde Peripherals
en ConfigTools con los siguientes parámetros:

\begin{table}[H]
    \centering
    \begin{tabular}{|c|c|c|}
        \hline
        \textbf{Campo}             & \textbf{Valor} & \textbf{Descripcion}                     \\
        \hline
        Enable DAC                 & Enabled        & Habilita el modulo DAC                   \\
        Reference voltage source   & VREF\_INT      & Fuente de voltaje de referencia          \\
        Opamp as buffer            & Enabled        & Habilita el buffer de salida             \\
        Enable LPDAC               & Enabled        & Habilita el modo de bajo consumo         \\
        Set LPDAC level after init & Enabled        & Configura el nivel de salida al encender \\
        LPDAC level                & 0              & Nivel de salida inicial del DAC (0V)     \\
        \hline
    \end{tabular}
    \caption{Configuración del DAC en ConfigTools/Peripherals}
\end{table}

\subsection{Programa Principal}

El programa principal inicializa los modulos (pines, clocks, perifericos y
consola de debugging) y habilita la interrupcion del CTIMER0. Luego ejecuta un
bucle infinito que actualiza el color de los LEDs RGB dependiendo de la
velocidad de muestreo seleccionada.

El modo de ejecucion del programa es el siguiente:

\begin{enumerate}
    \item Se presiona el switch de \texttt{SW3} para habilitar la adquisicion de muestras
          (\textbf{RUN}).
    \item Se presiona el switch de \texttt{SW2} para cambiar la velocidad de muestreo.
          (Por defecto la velocidad de muestreo es de \texttt{8 kS/s}). El cambio de
          velocidad de muestreo se indica con el color del LED RGB integrado de la placa.
    \item Las muestras adquiridas son almacenadas en un buffer circular de 512 muestras y
          enviadas al DAC al mismo tiempo para su visualizacion.
    \item Se presiona el switch de \texttt{SW1} para detener la adquisicion de muestras
          (\textbf{STOP}).
    \item Las muestras almacenadas en el buffer son enviadas al DAC para su visualizacion
          de forma continua siguiendo la velocidad de muestreo seleccionada.
\end{enumerate}

\subsubsection{Handlers de Interrupciones y Funciones}

Los handlers de interrupciones y funciones se utilizan para manejar los eventos
generador por los perifericos y botones de la placa y ejecutar las acciones
correspondientes.

\begin{enumerate}
    \item \texttt{CTIMER0\_IRQHandler():} Este handler se ejecuta cada vez que el
          timer hace match con el valor configurado. En este handler se eligen dos caminos,
          si la placa se encuentra en modo \textbf{START} se inicia la proxima conversion del ADC,
          si la placa se encuentra en modo \textbf{STOP} se envia al DAC la proxima muestra del buffer
          circular para su visualizacion.\@
    \item \texttt{ADC0\_IRQHandler():} Este handler se ejecuta cada vez que el
          ADC genera una interrupción por tener una nueva muestra en la FIFO.\@
          En este handler se lee la muestra de la FIFO, se almacena en el buffer
          circular y se envia al DAC para su visualizacion.
    \item \texttt{GPIO00\_IRQHandler():} Este handler se ejecuta cada vez que se presiona
          el switch de \texttt{SW3}. Habilita o deshabilita la adquisicion de
          muestras (\textbf{START} / \textbf{RUN}).
    \item \texttt{GPIO01\_IRQHandler():} Este handler se ejecuta cada vez que se presiona
          el switch de \texttt{SW2}. En este handler se cambia la velocidad de muestreo
          siguiendo un buffer circular de velocidades de muestreo.
    \item \texttt{setLedColors(rate):} Esta funcion se encarga de cambiar el color del LED RGB
          integrado de la placa dependiendo de la velocidad de muestreo seleccionada.
\end{enumerate}

La siguiente tabla muestra los colores del LED RGB integrados de la placa
dependiendo de la velocidad de muestreo seleccionada:

\begin{table}[H]
    \centering
    \begin{tabular}{|c|c|c|}
        \hline
        \textbf{Velocidad de muestreo} & \textbf{Color del LED} & \textbf{Configuracion RGB} \\
        \hline
        8 kS/s                         & \color{green}VERDE     & R = 100                    \\
        16 kS/s                        & \color{blue}AZUL       & B = 100                    \\
        22 kS/s                        & \color{yellow}AMARILLO & R = 100, G = 100           \\
        44 kS/s                        & \color{red}ROJO        & G = 100                    \\
        48 kS/s                        & \color{magenta}MAGENTA & R = 100, B = 100           \\
        \hline
    \end{tabular}
    \caption{Colores del LED RGB integrados de la placa}
\end{table}

\subsection{Resultados}

TODO

\end{document}